%%% Работа с русским языком
\usepackage{cmap}       % поиск в PDF
\usepackage{mathtext}      % русские буквы в формулах
\usepackage[T2A]{fontenc}   % кодировка
\usepackage[utf8]{inputenc}   % кодировка исходного текста
\usepackage[english,russian]{babel}   % локализация и переносы
\usepackage{setspace}
\usepackage{indentfirst}
%\usepackage{pbox}
\usepackage{pdfpages} %пакет для вставки pdf страниц

%%% Пакеты для работы с математикой
\usepackage{amsmath, amsfonts, amssymb, amsthm, mathtools}
\usepackage{icomma}
\numberwithin{equation}{section}

%% Номера формул
%\mathtoolsset{showonlyrefs=true} % Показывать номера только у тех формул, на которые есть \eqref{} в тексте.
%\usepackage{leqno}               % Нумерация формул слева
% Обновление нумерации уравнений
\renewcommand{\theequation}{\arabic{equation}}  % сквозная нумерация

%% Шрифты
\usepackage{euscript}   % Шрифт Евклид
\usepackage{mathrsfs}    % Красивый матшрифт

% Настройка стилей заголовков
\usepackage{titlesec} % Пакет для настройки заголовков
\titleformat{\section}
  {\normalfont\bfseries\centering} % Жирный, прописные
  {}                               % Нумерация (пусто - без номера)
  {0pt}                            % Отступ перед заголовком
  {}                               % Код перед заголовком

\titleformat{\subsection}
  {\normalfont\itshape}            % Курсив
  {}                               % Нумерация (пусто - без номера)
  {0pt}                            % Отступ перед заголовком
  {}                               % Код перед заголовком

%% Поля (геометрия страницы)
\usepackage[left=3cm,right=1.5cm,top=2cm,bottom=2cm,bindingoffset=0cm]{geometry}
\setlength{\topskip}{0pt}
%% Русские списки
\usepackage{enumitem}
\makeatletter
\AddEnumerateCounter{\asbuk}{\russian@alph}{щ}
\makeatother
\usepackage{pdflscape} % Для альбомной ориентации страницы
\usepackage{setspace}
\renewcommand{\baselinestretch}{0.8} % Уменьшение межстрочного интервала

\usepackage{fancyhdr}
\pagestyle{fancy}
\fancyhf{} % Очистить все заголовки и нижние колонтитулы
\fancyfoot[C]{\thepage} % Номер страницы в центре нижнего колонтитула
\renewcommand{\headrulewidth}{0pt} % Убрать линию, отделяющую верхний колонтитул от текста

%%% Работа с картинками
\usepackage{caption}
\captionsetup{format=plain, justification=justified, singlelinecheck=false}
\usepackage{longtable}
%\captionsetup{justification=centering} % центрирование подписей к картинкам
\usepackage{graphicx}                  % Для вставки рисунков
\graphicspath{{images/}{images2/}}     % папки с картинками
\setlength\fboxsep{3pt}                % Отступ рамки \fbox{} от рисунка
\setlength\fboxrule{1pt}               % Толщина линий рамки \fbox{}
\usepackage{wrapfig}                   % Обтекание рисунков и таблиц текстом
\usepackage{float} % для принудительного размещения рисунков

%%% Работа с таблицами
\usepackage{array,tabularx,tabulary,booktabs} % Дополнительная работа с таблицами
\usepackage{longtable}                        % Длинные таблицы
\usepackage{multirow}                         % Слияние строк в таблице
\usepackage{algorithm}
\usepackage{algpseudocode}
\usepackage{makecell} % Для многострочных заголовков в таблице


%% Красная строка
\setlength{\parindent}{2em}

%% Интервалы
\linespread{1.5}
\usepackage{multirow}

%% TikZ
\usepackage{tikz}
\usetikzlibrary{graphs,graphs.standard}

%% Верхний колонтитул
\usepackage{fancyhdr}
\pagestyle{fancy}
% Отключение отображения названия секции слева
\fancyhead[R]{}

%% Перенос знаков в формулах (по Львовскому)
\newcommand*{\hm}[1]{#1\nobreak\discretionary{}{\hbox{$\mathsurround=0pt #1$}}{}}


%% библиография
\frenchspacing
\usepackage[
    backend=biber,
    bibencoding=utf8,
    sorting=none,
    maxcitenames=1,
    maxbibnames=999,
    style=gost-numeric,
    language=auto,
    autolang=other,
    movenames=false]{biblatex}
\addbibresource{biblio.bib}
\usepackage[autostyle]{csquotes}

% блок исправляет (и др.) на (et al) для публикаций на латинице
\renewbibmacro*{name:andothers}{%
  \ifboolexpr{
    test {\ifnumequal{\value{listcount}}{\value{liststop}}}
    and
    test \ifmorenames
  }
    {\ifnumgreater{\value{liststop}}{1}
       {\finalandcomma}
       {}%
     \andothersdelim\bibstring{andothers}}
    {}}


%% дополнения
\usepackage{float}   % Добавляет возможность работы с командой [H] которая улучшает расположение на странице
%\usepackage{gensymb} % Красивые градусы
\usepackage{caption} % Пакет для подписей к рисункам, в частности, для работы caption*
\usepackage{tabto}

% подключаем hyperref и bookmark (для ссылок внутри  pdf)
\usepackage[
  colorlinks=false,
  unicode,
  pdftex]{hyperref}
\hypersetup{
  hidelinks, % Скрывает рамки и цвета ссылок
  }
%\usepackage[unicode, pdftex]{hyperref}
\usepackage{bookmark}


% для подсчета количества объектов в разделе "Структура и объем диссертации"
\usepackage{totcount}
\usepackage{totpages}
%\usepackage{cite}
%%% Объявление счетчиков для таблиц и рисунков
\regtotcounter{table}
\regtotcounter{figure}
%\regtotcounter{cite}